\Needspace{12\baselineskip}
\originalchapter{自由能、相稳定与粗粒化}
\originalsection{热力学积分}
平均能量不足以比较两个相的稳定性，因为构象熵可以逆转能量排序。令 $H_\lambda=(1-\lambda)H_0+\lambda H_1$，配分函数 $Z_\lambda$。对 $F_\lambda=-\kb T\ln Z_\lambda$ 求导有 $\partial_\lambda F=\avg{H_1-H_0}_\lambda$，因此
\[
F_1-F_0=\int_0^1\avg{H_1-H_0}_\lambda\dd\lambda.
\]
积分路径上必须能取得平衡样本；若穿过强一阶转变，直接插值可能遭遇滞后与采样困难，需要分段参考态或其他路径。数值积分误差与Monte Carlo采样误差是不同来源。

另一个恒等式 $F_1-F_0=-\kb T\ln\avg{\e^{-\beta(H_1-H_0)}}_0$ 看似只需一个端点，但由参考态中极少出现的构型主导时方差非常大。一个稳定的小误差条并不保证罕见权重已被采到。双向估计和相邻状态重叠检查更有意义。

\originalsection{模型Hamilton量和结构预测}
第一性原理能量可以拟合为占据变量的团簇展开：$E(\sigma)=J_0+\sum_\alpha J_\alpha\Phi_\alpha(\sigma)$，$\Phi_\alpha$ 是单点、配对或多点占据乘积的对称平均。它将昂贵的电子计算映射到可大规模采样的有效模型。训练能量少、拟合误差低仍可能漏掉相关构型，结构预测应保留独立验证和基态搜索的不完整性。

二元材料的零温稳定结构由形成能凸包决定，落在两结构连线之上的构型可通过分相降低能量。有限温度则需自由能凸包；化学势是共切线斜率。压力改变 $E+PV$，而温度改变熵贡献。不能据单个常压总能量差就确定所有高压、高温转变。

DFT给出的基态能对有限温度模型提供输入，但振动、电子和构型激发的自由能要按需要补上。高压比较尤其要求一致赝势及收敛，因为芯电子重叠或参考状态可迁移性可能发生改变。公开讲义中的行业和具体高压图是实例而非本卷运行结果，受限外图不重新分发。

\originalsection{Kinetic Monte Carlo与真实时间}
平衡Metropolis步通常没有物理时间。若系统处于离散亚稳态、转移事件率 $k_j$ 已知且等待近似指数分布，KMC总率 $K=\sum_jk_j$，下一事件以 $k_j/K$ 选择，等待时间 $\Delta t=-\ln u/K$。这两个随机抽样分别确定事件类型和时间，不能用每步固定一秒替代。

一维跳跃率 $k_+,k_-$、格距 $a$ 的精确平均漂移和扩散是 $v=a(k_+-k_-)$、$D=a^2(k_++k_-)/2$。代码实际运行事件序列，用独立路径验证这两个量。更复杂体系的率常由 $\nu\e^{-\beta\Delta E^\ddagger}$ 估计，前因子和势垒本身的不确定度必须保留。

粗粒化进一步消去快速或局部自由度，获得平均力势与有效输运。平衡分布匹配不保证时间尺度匹配，消去自由度还会产生记忆和噪声。第一性原理MD、加速MD和空间粗粒化的目的不同，本卷围绕这一区别建立方法联系，未把一份简单脚本说成全尺度材料预测程序。

\originalsection{自由能方法的推导、重叠和参考态}
$Z_\lambda=\int\e^{-\beta H_\lambda}\dd\Gamma$，对 $\lambda$ 微分得 $\partial_\lambda\ln Z=-\beta\avg{\partial_\lambda H}$，这证明热力学积分。若 $\lambda$ 改变质量、体积或可积变量测度，还要保留该测度的导数；不能一律只算势能差。

自由能微扰由 $Z_1/Z_0=\avg{\e^{-\beta\Delta H}}_0$ 得到。Jensen不等式给 $\Delta F\le\avg{\Delta H}_0$，反向又给 $\Delta F\ge\avg{\Delta H}_1$；两端差太大常提示分布重叠不足。少数大权重能支配指数平均，普通块误差可能看不见尚未采到的尾。把路径拆成相邻重叠态通常比一次跨越整个变化可靠。

简谐参考 $H_\lambda=k(\lambda)x^2/2$，$\avg{x^2}=\kb T/k$，积分给 $\Delta F=(\kb T/2)\ln(k_1/k_0)$。附带积分实际得到 $.693147194\kb T$，和 $k_1/k_0=4$ 的解析 $.693147181\kb T$ 一致。真实晶体Einstein参考还涉及中心质量约束与多个自由度，不能把一维例子的常数直接搬到完整材料熔点。

\originalsection{凸包、共切线和团簇展开的物理逻辑}
二元材料组成 $x$，若均匀构型自由能 $f(x)$ 高于两相组成 $x_1,x_2$ 的连线，按杠杆规则 $x=\lambda x_1+(1-\lambda)x_2$ 分相会更稳定。平衡端点满足 $f'(x_1)=f'(x_2)=[f(x_2)-f(x_1)]/(x_2-x_1)$，即两组分化学势一致。二阶导 $f''<0$ 的区间对组成涨落不稳定；它和两相共存区并不相同。

团簇展开将占据函数展开为正交基（例如二值 $s_i=\pm1$ 的乘积），晶格对称将等价团簇合并。拟合 $J_\alpha$ 时须留独立构型验证，避免高阶项把少量训练点完全插值却预测失真。实际基态搜索可能漏结构，有限温度MC又有采样误差，两个误差应分别报告。原讲义中的具体材料实例仅作为来源说明，未照搬受限图或编造第一性原理数据。

\originalsection{过渡态与KMC率的来源}
当体系在亚稳势阱内快速热化、跨势垒事件稀少，可用过渡态分割面上的平衡通量估计
\[k\simeq(\kb T/h)\exp(-\beta\Delta G^\ddagger)\]
再乘传输修正。经典谐振动近似的前因子也可由稳态与鞍点的稳定模频率比给出。势垒高度不是两端能量差，必须定位路径上的鞍点。沿反应坐标反复穿越使简单过渡态计数过大。

Nudged Elastic Band用若干路径图像、沿路径弹簧和垂直真实力定位最小能量路径；最高图像不必精确为鞍点，可用攀爬图像/鞍点搜索细化。初始路径可能通往不同机理，因此一种收敛路径不证明全球最低势垒。本文介绍方法和率的条件，未伪称运行真实材料的NEB。

已知若干事件率时，等待生存概率满足 $\dot P_0=-KP_0$，$P_0=\e^{-Kt}$，逆变换给 $\Delta t=-\ln u/K$，条件事件概率为 $k_j/K$。若率在等待期间随外场显著变化，必须积分时间依赖总率，固定 $K$ 的公式不再适用。记忆效应、刚进入势阱尚未热化也破坏Markov假设。

\originalsection{粗粒化和多尺度的边界}
设保留坐标 $Q(\Gamma)$，平均力势由 $\e^{-\beta F(Q)}\propto\int\delta[Q-Q(\Gamma)]\e^{-\beta H(\Gamma)}\dd\Gamma$ 定义。即使重现这个平衡分布，消去的变量仍可引入摩擦、噪声和记忆，简单Hamilton轨迹未必重现原始时间相关。粗粒化势通常随温度/组成而变，跨条件迁移需核验。

公开课程还讨论加速MD和更大尺度材料响应。加速稀有事件的方法要保留或校正物理时间；平衡偏置采样要重加权；连续本构模型又需有效参数与尺度分离。本卷在原子势、电子态、有限温度、稀有跃迁间建立完整方法链，代码分别验证可独立核验的小模型，明确哪些真实材料原软件任务没有实际重跑。

\originalsection{团簇函数的正交、完备与ECI投影}
3.320的4\_28公开讲义第6--8页给出团簇基。本段补足证明中的空团簇与计数约定。N个二元位点 $\sigma_i=\pm1$，对每个子集 $A\subseteq\{1,\ldots,N\}$ 定义 $\Phi_A(\sigma)=\prod_{i\in A}\sigma_i$，空集给 $\Phi_\varnothing=1$。以所有 $2^N$ 构型等权的内积 $\langle f,g\rangle=2^{-N}\sum_\sigma f(\sigma)g(\sigma)$，各自旋独立且均值零。

乘积 $\Phi_A\Phi_B=\Phi_{A\triangle B}$，因为交集中每个自旋平方等于1。若 $A\ne B$，对称差至少包含一个自旋，独立求和含该自旋的均值零；若 $A=B$，乘积恒为1。故 $\langle\Phi_A,\Phi_B\rangle=\delta_{AB}$。非空团簇均值零，\emph{空团簇的均值为1}，不能照简略幻灯片把所有基函数都说成零均值。

构型上实函数的线性空间维数为 $2^N$，而子集A也恰有 $2^N$ 个。已经正交归一的这组函数因此为完备基：任意能量函数可唯一展开
\[
 E(\sigma)=\sum_A J_A\Phi_A(\sigma),\qquad
 J_A=2^{-N}\sum_\sigma E(\sigma)\Phi_A(\sigma).
\]
证明不是仅靠把低阶项写出来，而是先给独立性与维数，再以内积提取系数。全基也允许任意高阶多体作用；实际截断到短程、低阶团簇是物理和数据拟合近似，不能同时宣称近似式对任意E精确。

若晶格对称性使同一轨道 $\mathcal O$ 的系数相同，可把等价团簇归组：$E=\sum_{\mathcal O}M_{\mathcal O}J_{\mathcal O}\bar\Phi_{\mathcal O}$，其中 $M_{\mathcal O}$ 是该轨道的团簇数、$\bar\Phi$ 为轨道平均；每原子能量还须除以N。平均后的函数不再沿用单个基函数的单位范数。若改用有序(i,j,k)求和，成对、三体项分别除以 $2!$、$3!$ 消去排列重复，和本卷“无序对一次”的Ising约定必须对应。

例如两个交换等价位点，$E=J_0+J_1(\sigma_1+\sigma_2)+J_2\sigma_1\sigma_2$。用四个构型能量可求 $J_0=(E_{++}+E_{+-}+E_{-+}+E_{--})/4$、$J_1=(E_{++}-E_{--})/4$、$J_2=(E_{++}+E_{--}-E_{+-}-E_{-+})/4$。给定组成的构型子集没有遍历所有自旋，等权正交证明不再直接成立；实际DFT训练集应构造设计矩阵、检查秩、留出结构验证，防止高阶拟合在少量样本上插值后错误外推。ECI一旦确定才进入MC采样，拟合误差与采样误差是两个来源。
